% Generated by geometry_pipeline.latex_renderer. Do not edit by hand.

\chapter*{Вступні поняття}

\addcontentsline{toc}{chapter}{Вступні поняття}

\begin{tcolorbox}[colback=boxbglight,colframe=primaryblue,title={ЯК ЧИТАТИ ЦЕЙ ВСТУП},fonttitle=\bfseries\sffamily,breakable]

Параграфи 1--12 вводять базову мову геометрії. Основний текст є сучасною українською адаптацією Кисельова; редакторські уточнення виділено окремо.

\end{tcolorbox}

% block: geo.p009.s01.geometric-figures; printed source pages: 9

\paragraph{1. Геометричні фігури.}

Частину простору, яку займає фізичне тіло, називають \textbf{геометричним тілом}.

Геометричне тіло відокремлюється від навколишнього простору \textbf{поверхнею}. Частина поверхні відокремлюється від сусідньої частини \textbf{лінією}, а частина лінії — \textbf{точкою}.

Геометричне тіло, поверхня, лінія і точка не існують окремо. Проте за допомогою абстрагування ми можемо розглядати поверхню незалежно від тіла, лінію — незалежно від поверхні, а точку — незалежно від лінії. У геометричній моделі поверхня не має товщини, лінія не має ні товщини, ні ширини, а точка не має розмірів.

Будь-яку сукупність точок, ліній, поверхонь або тіл, розміщених у просторі певним чином, називають \textbf{геометричною фігурою}.

% block: geo.p009.s02.geometry; printed source pages: 9

\paragraph{2. Геометрія.}

Науку про властивості геометричних фігур називають \textbf{геометрією}. Назва походить від грецьких слів, пов'язаних із вимірюванням землі: у давнину одним із головних практичних завдань геометрії було вимірювання відстаней і площ на земній поверхні.

% block: geo.p009.s03.straight-line; printed source pages: 9

\paragraph{3. Пряма.}

Наочне уявлення про пряму дає тонка натягнута нитка. Пряма має основну властивість: \textbf{через будь-які дві різні точки можна провести одну й тільки одну пряму}.

Звідси випливає: якщо дві прямі мають дві спільні точки, то вони збігаються. Тому дві різні прямі можуть перетинатися не більш ніж в одній точці.

\motionlink{Як рух точки народжує лінію}%
{Порівняй рух зі сталим і змінним напрямом та простеж, коли слід точки є прямим, а коли стає кривим.}%
{https://radchr.github.io/Algebra/geometry-point-trace.html}%
{qr_point_trace.png}

% block: geo.p009.s04.line-segment-ray; printed source pages: 9, 10

\paragraph{4. Пряма, відрізок і промінь.}

Пряма нескінченно продовжується в обидва боки (рис. 1). Її позначають малою літерою, наприклад $a$, або двома великими літерами, поставленими біля двох її точок: пряма $AB$.

Частину прямої, обмежену з обох боків двома точками, називають \textbf{відрізком} (відрізок $CD$, рис. 2). Ці точки називають кінцями відрізка.

Частину прямої, обмежену з одного боку початковою точкою, називають \textbf{променем} (рис. 3). Промінь $AB$ починається в точці $A$ і проходить через точку $B$.

\begin{center}
\begin{tikzpicture}[scale=1.0,>=Stealth]
  \draw[dashed] (-3.4,0) -- (-2.7,0);
  \draw[<->,very thick,color=primaryblue] (-2.7,0) -- (2.7,0);
  \draw[dashed] (2.7,0) -- (3.4,0);
  \fill (-1.6,0) circle (1.8pt) node[above] {$A$};
  \fill (1.7,0) circle (1.8pt) node[above] {$B$};
  \node[below] at (0,0) {$a$};
\end{tikzpicture}

{\small\textbf{Рис. 1. Пряма AB, позначена також малою літерою a}}\\[6pt]
\begin{tikzpicture}[scale=1.0]
  \draw[very thick,color=primaryblue] (-2.2,0) -- (2.2,0);
  \fill (-2.2,0) circle (2pt) node[above] {$C$};
  \fill (2.2,0) circle (2pt) node[above] {$D$};
\end{tikzpicture}

{\small\textbf{Рис. 2. Відрізок CD}}\\[6pt]
\begin{tikzpicture}[scale=1.0,>=Stealth]
  \fill (-2.2,0) circle (2pt) node[above] {$A$};
  \draw[->,very thick,color=primaryblue] (-2.2,0) -- (2.8,0);
\end{tikzpicture}

{\small\textbf{Рис. 3. Промінь із початком у точці A}}\\[6pt]
\end{center}

\begin{feynmanbox}[Сучасні назви]
Кисельов уживав вислови «кінцева пряма» та «півпряма». У сучасній українській шкільній геометрії використовуємо терміни \textbf{відрізок} і \textbf{промінь}.
\end{feynmanbox}

% block: geo.p010.s05.segment-equality; printed source pages: 10

\paragraph{5. Рівність і нерівність відрізків.}

Два відрізки вважають рівними, якщо їх можна накласти один на один так, щоб збіглися обидва кінці. Наприклад, відрізок $AB$ накладають на $CD$ (рис. 4), суміщаючи $A$ з $C$. Якщо при цьому $B$ збігається з $D$, то $AB=CD$. Якщо ні, меншим вважають той відрізок, який становить частину іншого.

Щоб відкласти на прямій відрізок, рівний даному, використовують циркуль.

\begin{center}
\begin{tikzpicture}[scale=0.95]
  \draw[very thick,color=primaryblue] (-3.2,0.5) -- (-0.4,0.5);
  \fill (-3.2,0.5) circle (2pt) node[above] {$A$};
  \fill (-0.4,0.5) circle (2pt) node[above] {$B$};
  \draw[very thick,color=feynmangreen] (0.6,0.5) -- (3.4,0.5);
  \fill (0.6,0.5) circle (2pt) node[above] {$C$};
  \fill (3.4,0.5) circle (2pt) node[above] {$D$};
  \draw[dashed] (-3.2,-0.35) -- (-0.4,-0.35);
  \draw[dashed] (0.6,-0.35) -- (3.4,-0.35);
  \node at (0.1,-0.35) {$AB=CD$};
\end{tikzpicture}

{\small\textbf{Рис. 4. Порівняння відрізків AB і CD}}\\[6pt]
\end{center}

% block: geo.p010.s06.segment-sum; printed source pages: 10

\paragraph{6. Сума відрізків.}

Суму відрізків $AB$, $CD$, $EF$ будують так (рис. 5). На прямій беруть точку $M$ і послідовно відкладають в одному напрямку $MN=AB$, $NP=CD$, $PQ=EF$. Тоді:
\[
MQ=AB+CD+EF.
\]

Сума відрізків не залежить від порядку доданків і від способу їх групування — це переставна та сполучна властивості додавання.

\begin{center}
\begin{tikzpicture}[scale=0.9]
  \draw[very thick] (-4,1.2) -- (-2.2,1.2) node[midway,above] {$AB$};
  \draw[very thick] (-1.4,1.2) -- (0.2,1.2) node[midway,above] {$CD$};
  \draw[very thick] (1.0,1.2) -- (2.4,1.2) node[midway,above] {$EF$};
  \draw[very thick,color=primaryblue] (-4,-0.2) -- (2.4,-0.2);
  \foreach \x/\name in {-4/M,-2.2/N,-0.6/P,0.8/Q} {
    \fill (\x,-0.2) circle (2pt) node[below] {$\name$};
  }
  \node[above] at (-3.1,-0.2) {$AB$};
  \node[above] at (-1.4,-0.2) {$CD$};
  \node[above] at (0.1,-0.2) {$EF$};
\end{tikzpicture}

{\small\textbf{Рис. 5. Послідовне додавання трьох відрізків}}\\[6pt]
\end{center}

% block: geo.p010.s07.segment-operations; printed source pages: 10, 11

\paragraph{7. Дії над відрізками.}

Різницею відрізків $AB$ і $CD$, де $AB>CD$, називають такий відрізок, який разом із $CD$ дає $AB$. Добуток довжини відрізка на натуральне число означає суму відповідної кількості рівних відрізків; ділення на натуральне число — поділ відрізка на рівні частини.

Коли відрізки виміряно в однакових одиницях, геометричним діям відповідають дії над їхніми числовими довжинами.

% block: geo.p011.s08.circle-concepts; printed source pages: 11

\paragraph{8. Коло і пов'язані поняття.}

Якщо одну ніжку циркуля поставити в точку $O$, а другою обертатися навколо неї, олівець опише \textbf{коло} — замкнену лінію, усі точки якої однаково віддалені від центра $O$ (рис. 6).

Відрізки $OA$, $OB$, $OC$, що сполучають центр із точками кола, називають \textbf{радіусами}. Усі радіуси одного кола рівні. Пряму, що перетинає коло у двох точках, називають \textbf{січною}; відрізок із кінцями на колі — \textbf{хордою}; хорду, що проходить через центр, — \textbf{діаметром}; частину кола — \textbf{дугою}.

Частину площини, обмежену колом, називають \textbf{кругом}. Частина круга між двома радіусами є \textbf{сектором}, а частина круга між хордою та відповідною дугою — \textbf{сегментом}.

\begin{center}
\begin{tikzpicture}[scale=1.0]
  \coordinate (O) at (0,0);
  \draw[thick] (O) circle (2.2);
  \coordinate (A) at (0:2.2);
  \coordinate (B) at (48:2.2);
  \coordinate (C) at (90:2.2);
  \coordinate (D) at (180:2.2);
  \coordinate (E) at (235:2.2);
  \coordinate (F) at (340:2.2);
  \fill[feynmangreen,opacity=.16] (O) -- (A) arc (0:48:2.2) -- cycle;
  \fill[mediapurple,opacity=.14] (E) arc (235:340:2.2) -- cycle;
  \draw[thick] (D) -- (A);
  \draw[thick] (O) -- (B);
  \draw[thick] (O) -- (C);
  \draw[thick] (E) -- (F);
  \draw[dashed] ($(E)!-.35!(F)$) -- ($(E)!1.35!(F)$);
  \fill (O) circle (2pt) node[below left] {$O$};
  \foreach \p in {A,B,C,D,E,F} {\fill (\p) circle (1.6pt) node[above right] {$\p$};}
  \node at (1.1,0.55) {сектор};
  \node at (0,-1.75) {сегмент};
\end{tikzpicture}

{\small\textbf{Рис. 6. Основні елементи кола і круга}}\\[6pt]
\end{center}

\begin{warningbox}[Коло і круг — різні об'єкти]
\textbf{Коло} — це лише межа, тобто замкнена лінія. \textbf{Круг} містить і межу, і всю внутрішню область. Англійською це відповідає словам \emph{circle} і \emph{disk}.
\end{warningbox}

\motionlink{Коло як слід точки}%
{Рухай точку навколо нерухомого центра й спостерігай, як стала відстань до центра породжує коло.}%
{https://radchr.github.io/Algebra/geometry-circle-trace.html}%
{qr_circle_trace.png}

% block: geo.p011.s09.arc-equality; printed source pages: 11

\paragraph{9. Рівність і нерівність дуг.}

Дві дуги одного кола або рівних кіл вважають рівними, якщо вони суміщаються накладанням. Якщо одна дуга становить частину іншої, то вона менша.

% block: geo.p011.s10.arc-sum; printed source pages: 11, 12

\paragraph{10. Сума дуг.}

Суму дуг однакового радіуса будують послідовним відкладанням на тому самому колі. На рис. 7 від точки $M$ відкладають дугу $MN$, рівну першій даній дузі, а потім дугу $NP$, рівну другій. Тоді дуга $MP$ є їхньою сумою.

Сума дуг має переставну та сполучну властивості. Аналогічно визначають різницю дуг, множення дуги на число та ділення дуги на число.

\begin{center}
\begin{tikzpicture}[scale=1.0]
  \draw[thick] (0,0) circle (2.2);
  \coordinate (M) at (170:2.2);
  \coordinate (N) at (105:2.2);
  \coordinate (P) at (35:2.2);
  \draw[very thick,color=primaryblue] (M) arc (170:105:2.2);
  \draw[very thick,color=feynmangreen] (N) arc (105:35:2.2);
  \foreach \p/\name in {M/M,N/N,P/P} {
    \fill (\p) circle (2pt) node[above] {$\name$};
  }
  \fill (0,0) circle (2pt) node[below] {$O$};
  \node at (0,2.65) {$\overset{\frown}{MP}=\overset{\frown}{MN}+\overset{\frown}{NP}$};
\end{tikzpicture}

{\small\textbf{Рис. 7. Додавання дуг однакового радіуса}}\\[6pt]
\end{center}

% block: geo.p012.s11.plane; printed source pages: 12

\paragraph{11. Площина.}

Наочне уявлення про площину дає поверхня рівного скла або спокійної води. Для площини важливі дві властивості:

\begin{enumerate}
\item Якщо дві точки прямої лежать у площині, то вся пряма лежить у цій площині.
\item Будь-яку частину площини можна сумістити з іншим місцем тієї самої або іншої площини, зокрема після перевертання.
\end{enumerate}

% block: geo.p012.s12.geometry-branches; printed source pages: 12

\paragraph{12. Поділ геометрії.}

Геометрію поділяють на \textbf{планіметрію} і \textbf{стереометрію}. Планіметрія вивчає властивості фігур, усі точки яких лежать в одній площині. Стереометрія вивчає властивості просторових фігур, які не можна цілком розмістити в одній площині.
